Most published agentic security systems converge on a single execution loop. PentestGPT \cite{deng2024pentestgpt} drives a remote shell through a planner; CAI \cite{mayoral2025cai} adds a constrained Python tool registry and a $\$10$ per-run hard cap; Ji et al.\ \cite{ctfagent2025} measure and augment LLMs for CTF solving at CCS~2025; 
%RedTeamLLM \cite{redteamllm2025} adds plan-correction and bounded memory and surpasses PentestGPT on $11/13$ VMs. 
Across this body of work, the scaffold itself is treated as an implementation detail of a single best agent. The complementarity question we pose, \emph{do different scaffolds fail on different challenges}, is rarely instrumented because the underlying experiments do not run multiple scaffolds against the same items.

G-CTR \cite{mayoralvilches2025gametheoretic} demonstrated that injecting a Nash-equilibrium digest of the agent's own attack graph back into its context lifts cyber-range success from 20.0\% to 42.9\% with a 2.7$\times$ better cost-per-success and a 5.2$\times$ reduction in behavioural variance. The present paper adopts the same bio-inspired\todo{[review C5.5] "bio-inspired" appears here, in \S5.4 (multiagent.tex line 75) and \S7.3 (discussion.tex line 21); the analogy is to human expert teams (sociological), not biological mechanisms. Replace with "team-of-experts-inspired" or reuse the "cognitive-heterogeneity posture" wording already used in the introduction; one global substitution covers all three sites.} posture, namely inserting a human-style heuristic into the loop, but moves from \emph{game theory} as the heuristic to \emph{cognitive heterogeneity}, and from \emph{intra-agent guidance} to \emph{inter-agent coordination}.

Beyond cybersecurity, the value of diverse reasoning is well established. Self-consistency \cite{wang2023selfconsistency} samples multiple chains of thought from a single model and marginalises over answers, gaining $+17.9$ percentage points on GSM8K. AlphaCode \cite{li2022alphacode} generates up to $10^6$ candidate programmes per task and clusters them by behavioural equivalence. Recent ranked-vote self-consistency \cite{rankedvoting2025}\todo{[review C2.1] bib entry rankedvoting2025 has author={Anonymous}; verify against arXiv:2505.10772 and replace with named authors.} refines this further by weighting answers within each sample. These results establish a pattern: the union of diverse reasoning paths exceeds the Pass@$K$ of any individual path. We test the same pattern at a coarser granularity, the scaffold rather than the sample.

Multi-agent LLM frameworks, CAMEL \cite{li2023camel}, AutoGen \cite{wu2024autogen}, MetaGPT \cite{hong2024metagpt}, ChatDev \cite{qian2024chatdev}, and AgentVerse \cite{chen2024agentverse}, demonstrate that role-specialised LLM agents communicating in natural language can outperform a single agent on tasks from software development to scientific reasoning. MacNet \cite{macnet2025} scales to $> 1000$ agents over a directed acyclic graph, validating that peer-to-peer substrates outperform coordinator bottlenecks at scale. The blackboard pattern, formalised in Hearsay-II \cite{hearsay1980} and developed by Hayes-Roth \cite{hayesroth1985blackboard}, decouples specialist knowledge sources from a shared workspace. Recent work has revived it for LLM agents \cite{blackboardllm2025,blackboarddatascience2024}\todo{[review C2.1] bib entry blackboarddatascience2024 has author={Anonymous} and a year/arXiv-ID mismatch (year=2024 but arXiv ID 2510.01285 implies Oct 2025); verify and replace.}; we adopt this pattern as the integration substrate for our scaffold-heterogeneous architecture (Section~\ref{sec:multiagent}).

In offensive security specifically, Co-RedTeam \cite{coredteam2026}\todo{[review C2.2] coredteam2026 bib entry has author={He, Yifeng and others} and arXiv ID 2602.02164; this is the single counter-anchor for the novelty claim below, expand the author list and verify the ID resolves.} orchestrates discovery and exploitation as separate agents that exchange validated findings. \todo{[review C5.2] this is a three-conjunct "to our knowledge" first-to claim with a single counter-anchor; document the search scope (databases, queries, date) in a footnote, or merge below, suggestion: "we are not aware of prior work that does all three of \{scaffold heterogeneity, per-item agreement instrumentation, derived multi-agent architecture\}".} To our knowledge, no prior work runs multiple structurally different scaffolds (rather than role-specialised instances of the same scaffold) on the same target, instruments their per-item agreement, and uses the empirical complementarity to derive a multi-agent architecture.
